% !TEX program = xelatex
%==============================================================================
% Muhammad Shaheer Luqman — Resume (research-focused, single column)
% Theme matched to portfolio: accent blue #3f63ad, heading navy #1a2744,
% body slate #5a6a8a. Font: Poppins.
% This variant leads with Research & Publications and is intentionally lighter
% and less dense than the AI Engineer variant.
% IMPORTANT: Compile with xelatex (uses local Poppins font via fontspec):
%   xelatex Shaheer_Resume.tex
%==============================================================================
\documentclass[10pt,a4paper]{article}

\usepackage{amssymb}
\usepackage{fontspec}
% Poppins — loaded by file from the local fonts/ folder so the resume compiles
% on any machine without a system-wide font install. Geometric sans, slightly
% heavier than Arial.
\setmainfont{Poppins}[
  Path = fonts/,
  Extension = .ttf,
  UprightFont = Poppins-Light,
  ItalicFont  = Poppins-LightItalic,
  BoldFont    = Poppins-SemiBold,
  BoldItalicFont = Poppins-SemiBoldItalic,
  Scale = 0.90
]
% Slightly more generous margins than the dense variant so the page breathes.
\usepackage[top=1.2cm,bottom=1.2cm,left=1.6cm,right=1.6cm]{geometry}
\usepackage[dvipsnames,table]{xcolor}
\usepackage{enumitem}
\usepackage{titlesec}
\usepackage{fontawesome5}
\usepackage{hyperref}
\usepackage{tabularx}
\usepackage{array}
\usepackage{parskip}
\usepackage{etoolbox}
\usepackage{multicol}
\setlength{\multicolsep}{2pt}
\setlength{\columnsep}{1.2em}
% Roomier paragraph spacing than the dense variant.
\setlength{\parskip}{3pt plus 1pt minus 1pt}

% Consistent inter-line spacing across the whole document.
\usepackage{setspace}
\setstretch{1.02}

%------------------------------- Colors --------------------------------------
\definecolor{accent}{HTML}{3F63AD}   % portfolio primary blue
\definecolor{accentdk}{HTML}{2E4A84} % darker blue (hover state)
\definecolor{heading}{HTML}{1A2744}  % navy headings
\definecolor{body}{HTML}{4A5878}     % slate body text
\definecolor{muted}{HTML}{7A89A8}    % muted / dates
\definecolor{rule}{HTML}{3F63AD}
\definecolor{chipbg}{HTML}{E8EEF7}   % tech-tag background

%------------------------------- Hyperlinks ----------------------------------
\hypersetup{
  colorlinks=true,
  urlcolor=accent,
  linkcolor=accent,
  pdftitle={Muhammad Shaheer Luqman — Resume},
  pdfauthor={Muhammad Shaheer Luqman}
}
% Underline + italicize hyperlink text (in addition to the accent color).
\usepackage[normalem]{ulem}
\let\orighref\href
\renewcommand{\href}[2]{\orighref{#1}{\uline{\textit{#2}}}}

%------------------------------- Section style -------------------------------
\titleformat{\section}
  {\Large\bfseries\color{heading}}
  {}{0em}
  {}[\vspace{2pt}{\color{rule}\titlerule[1.2pt]}]
\titlespacing*{\section}{0pt}{10pt}{7pt}

% Default body color
\AtBeginDocument{\color{body}}

%------------------------------- Helper macros -------------------------------
% Experience / entry header: title (left, bold navy) + date (right, muted)
% #1 = title, #2 = date, #3 = subtitle (accent italic), #4 = meta (muted small)
% Pass empty {} for any unused field.
\newlength{\datew}
\newcommand{\entryline}[2]{%
  \noindent
  \begingroup
  \settowidth{\datew}{\color{muted}#2}%
  \parbox[t]{\dimexpr\linewidth-\datew-1em\relax}{\raggedright\bfseries\color{heading}#1}%
  \hfill
  \parbox[t]{\datew}{\raggedleft\color{muted}#2}%
  \par
  \endgroup
}
\newcommand{\subline}[1]{\ifblank{#1}{}{{\color{accent}\itshape #1}\par}}
\newcommand{\metaline}[1]{\ifblank{#1}{}{{\color{muted}\small #1}\par}}

% Publication entry: single row, two columns. Left column = paper title
% (wraps freely). Right column = status/venue (e.g. "Under Review",
% "ICETST 2024") stacked directly above the journal/arXiv reference, both
% right-aligned within the same box so they read as one unit. Colors match
% the Education/Experience \gridentry pattern: top line accent blue (like
% location), bottom line muted gray (like date).
% #1 = title, #2 = status/venue, #3 = journal/arXiv reference.
\newlength{\pubw}
\newcommand{\pubentryline}[3]{%
  \noindent\begingroup
  \settowidth{\datew}{\itshape#2}%
  \settowidth{\pubw}{\itshape\small #3}%
  \ifdim\pubw>\datew \setlength{\datew}{\pubw}\fi
  \parbox[t]{\dimexpr\linewidth-\datew-1em\relax}{\raggedright\bfseries\color{heading}#1}%
  \hfill
  \parbox[t]{\datew}{\raggedleft\itshape\color{accent}#2\\[2pt]{\itshape\small\color{muted}#3}}%
  \par
  \endgroup
}
\newcommand{\pubentry}[3]{%
  \par\vspace{-\parskip}\vspace{\entryskip}\vspace{-2pt}%
  \pubentryline{#1}{#2}{#3}%
}
% Consistent gap before every entry across all sections.
\newlength{\entryskip}\setlength{\entryskip}{10pt}
\newcommand{\entry}[4]{%
  \par\vspace{-\parskip}\vspace{\entryskip}\vspace{-2pt}%
  \entryline{#1}{#2}%
  \subline{#3}%
  \metaline{#4}%
}
% Use right after \section to cancel the leading gap on the first entry, so
% the heading-to-first-item spacing comes only from \titlespacing's "after"
% value, not from \entryskip (which controls spacing between later entries).
\newcommand{\firstentry}{\vspace{\parskip}\vspace{-\entryskip}\vspace{2pt}}

% Two-line entry header (Education / Experience): single row, two columns,
% each holding two stacked lines joined with the same explicit \\[2pt] gap
% used by \pubentryline, so the line-to-line spacing matches exactly.
% Left column: institution/company (bold navy) above degree/position (muted
% small caps). Right column: location (italic accent) above date range
% (italic muted small). #1 = name, #2 = location, #3 = degree/position,
% #4 = date range.
\newlength{\entrylocw}
\newcommand{\gridentryline}[4]{%
  \noindent\begingroup
  \settowidth{\entrylocw}{\itshape #2}%
  \settowidth{\datew}{\itshape\small #4}%
  \ifdim\datew>\entrylocw \setlength{\entrylocw}{\datew}\fi
  \parbox[t]{\dimexpr\linewidth-\entrylocw-1em\relax}{\raggedright\bfseries\color{heading}#1\\[2pt]{\small\color{muted}\MakeUppercase{#3}}}%
  \hfill
  \parbox[t]{\entrylocw}{\raggedleft\itshape\color{accent}#2\\[2pt]{\itshape\small\color{muted}#4}}%
  \par
  \endgroup
}
\newcommand{\gridentry}[4]{%
  \par\vspace{-\parskip}\vspace{\entryskip}\vspace{-2pt}%
  \gridentryline{#1}{#2}{#3}{#4}%
}

% Bullet list tuned tight
\newlist{tight}{itemize}{1}
\setlist[tight]{leftmargin=1.2em,itemsep=4pt,topsep=2pt,parsep=0pt,partopsep=0pt,
  label={\color{accent}\raisebox{0ex}{\normalsize\textbullet}}}

% Tech / skill chip
\newcommand{\chip}[1]{%
  \setlength{\fboxsep}{2pt}%
  \colorbox{chipbg}{\color{accentdk}\small\strut\,#1\,}\hspace{2pt}}

\setlength{\parindent}{0pt}

% Remove page numbers.
\pagestyle{empty}

%==============================================================================
\begin{document}

%------------------------------- Header --------------------------------------
\begin{center}
{\fontsize{26}{30}\selectfont\bfseries\color{accent}Muhammad Shaheer Luqman}\\[5pt]
{\footnotesize\color{heading}
\faEnvelope[regular]\,\href{mailto:muhammadshaheerluqman@gmail.com}{muhammadshaheerluqman@gmail.com}
~\textbar~\faHome\,\href{https://shaheerluqman.github.io/}{shaheerluqman.github.io}
~\textbar~\faGithub\,\href{https://github.com/ShaheerLuqman}{ShaheerLuqman}
~\textbar~\faLinkedin\,\href{http://linkedin.com/in/shaheerluqman}{shaheerluqman}
~\textbar~\faGraduationCap\,\href{https://scholar.google.com/citations?user=I3TSLaAAAAAJ&hl=en}{M. S. Luqman}
}
\end{center}

\vspace{3pt}
{\color{body}AI researcher working in reinforcement learning, computer vision, and robot manipulation. Published across arXiv (2$\times$) and IEEE, ICPC Gold Medalist, and summa cum laude graduate (3.90/4.00).}

%------------------------------- Education -----------------------------------
\section{Education}
\firstentry
\gridentry{National University of Computer and Emerging Sciences (FAST)}{Karachi, Pakistan}{BS Computer Science}{Sep 2021 -- Jun 2025}
\begin{tight}
  \item CGPA: \textbf{3.90 / 4.00} (Ranked \textbf{9 / 440})
  \item Summa Cum Laude Honors
\end{tight}

%------------------------------- Research Experience -------------------------
\section{Experience}
\firstentry
\gridentry{Retrocausal}{Remote (US-based)}{Applied Research Engineer}{Jun 2025 -- Present}
\vspace{3pt}
\begin{tight}
  \item Research on robot manipulation and Vision-Language-Action systems, unsupervised action segmentation, and temporal long-video understanding --- \textbf{two papers submitted} (see \textit{Research \& Publications}).
  \item Built and shipped end-to-end AI systems for industrial clients (incl.\ \textbf{BMW}), from model training to GPU deployment on \textbf{AWS} and \textbf{RunPod}.
\end{tight}
\vspace{3pt}
{\footnotesize\chip{Reinforcement Learning}\chip{Robotics}\chip{VLMs}\chip{Computer Vision}\chip{PyTorch}}

\gridentry{Retrocausal}{Remote (US-based)}{Software Engineer Intern}{Mar 2024 -- Jun 2024}
\vspace{3pt}
\begin{tight}
  \item Built a video activity-recognition pipeline with \textbf{body pose estimation} and \textbf{vision-language models} for automated industrial workflow analysis.
  \item Developed \textbf{LLM prompt-engineering} systems for analysis of industrial process data.
\end{tight}
\vspace{3pt}
{\footnotesize\chip{Computer Vision}\chip{VLMs}\chip{LLMs}\chip{MediaPipe}\chip{Python}}

%------------------------- Research & Publications ---------------------------
\section{Research \& Publications}
\firstentry
\pubentry{Enhancing Human-Likeness in Reinforcement Learning Agents via Hierarchical Macro Action Quantization}{Under Review}{arXiv:2605.30928}
\vspace{3pt}

A human-like RL framework (HiMAQ) that encodes human demonstrations into a two-level hierarchy of
macro actions for dexterous humanoid-hand control. On D4RL it improves human-likeness over the
non-hierarchical baseline while matching or exceeding success rates across IQL, SAC, and RLPD.
\href{https://arxiv.org/abs/2605.30928}{View on arXiv}
\vspace{3pt}

{\footnotesize\chip{Robot Manipulation}\chip{Hierarchical Action Quantization}\chip{Reinforcement Learning}\chip{Imitation Learning}}

\pubentry{Unsupervised Skeleton-Based Action Segmentation via Hierarchical Spatiotemporal Vector Quantization}{Under Review}{arXiv:2604.15196}
\vspace{3pt}

A hierarchical spatiotemporal vector-quantization framework for unsupervised skeleton-based action
segmentation, achieving state-of-the-art on the HuGaDB, LARa, and BABEL benchmarks.
\href{https://arxiv.org/abs/2604.15196}{View on arXiv}
\vspace{3pt}

{\footnotesize\chip{Unsupervised Learning}\chip{Action Segmentation}\chip{Vector Quantization}\chip{Skeleton-Based Recognition}}

\pubentry{Ensuring Child Safety: An IoT-Based Surveillance System for Remote Monitoring and Detection of Anomalous Behavior}{ICETST 2024}{IEEE Xplore}
\vspace{3pt}

An IoT surveillance system pairing hardware sensors with AI-driven anomaly detection for real-time
child-safety alerts. \href{https://ieeexplore.ieee.org/document/10737980}{View on IEEE Xplore}
\vspace{3pt}

{\footnotesize\chip{IoT Systems}\chip{Anomaly Detection}\chip{Surveillance}\chip{Embedded Systems}}

\vspace{10pt}
{\color{heading}\textbf{Research Interests:}}\ {\color{body}Robot Manipulation, Computer Vision, Reinforcement Learning, Vision-Language Models, IoT + AI Systems}

%------------------------------- Certifications -------------------------------
\section{Certifications}
\firstentry
\entry{IELTS Academic \textnormal{\small— English Proficiency}}{Jul 2026}{}{Overall Band: \textbf{7.5 / 9} \textbullet\ Reading \textbf{9.0} \textbullet\ Listening \textbf{8.5}}

%--------------------------- Technical Skills ---------------------------------
\section{Technical Skills}

\newcommand{\skillrow}[2]{\noindent{\color{heading}\textbf{#1}}\hspace{4pt}{\color{body}#2}\par\vspace{4pt}}

\skillrow{Languages}{Python, C / C++, JavaScript.}
\skillrow{ML Frameworks}{PyTorch, Hugging Face, Scikit-learn; YOLO, MediaPipe; Stable Diffusion; BERT, Whisper.}
\skillrow{Software Frameworks}{Django, FastAPI, Flask; Express.js, React, Next.js, Vue.}
\skillrow{Developer Tools}{AWS, RunPod; Docker, Git, Linux.}

%------------------------------- Selected Projects ----------------------------
\section{Projects}
\firstentry
\entry{PixelForge — Generative AI Marketing Platform}{}{Final Year Project (\href{https://github.com/ShaheerLuqman/PixelForge}{GitHub})}{}
Full-stack generative pipeline combining \textbf{BLIP} captioning, \textbf{Gemini} text generation, and
diffusion-based image and video synthesis for product-aware marketing content.

\entry{LushWear — Business Operation Management}{}{Client (\href{https://github.com/ShaheerLuqman/LushWear-IMS}{GitHub})}{}
\textbf{Production-ready system, built, sold, and actively maintained} for a local business to run its
day-to-day operations. Full-stack inventory, order, and storefront management with a \textbf{FastAPI}
backend, \textbf{Electron} desktop app, \textbf{Supabase} real-time database, and two-way \textbf{Shopify}
integration.

\vspace{2pt}
{\color{body}More on my \href{https://shaheerluqman.github.io/}{portfolio} and \href{https://github.com/ShaheerLuqman}{GitHub}.}

%--------------------------- Honors & Awards ---------------------------------
\section{Honors \& Awards}
\begin{tight}
  \item \textbf{ICPC Gold Medalist} — ICPC National Competitive Programming \textit{(Nov 2024)}
  \item \textbf{Summa Cum Laude} — Graduated with a high GPA (3.90 / 4.00), National University of Computer and Emerging Sciences (FAST) \textit{(2025)}
  \item \textbf{ICPC Bronze Medalist} — ICPC National Competitive Programming \textit{(Feb 2024)}
  \item \textbf{Winner} — Code Sprint, PROCOM 2023 \textit{(2023)}
  \item \textbf{Best Team Award} — Coder's Cup, ACM FAST
  \item \textbf{4x Rector's List (4.0 GPA)} — NUCES (FAST), 4 semesters
  \item \textbf{4x Dean's List (3.5+ GPA)} — NUCES (FAST), 4 semesters
\end{tight}

\end{document}
